\documentclass[11pt,a4paper]{article}
\usepackage[utf8]{inputenc}
\usepackage[T1]{fontenc}
\usepackage{geometry}
\geometry{a4paper, margin=2.5cm}
\usepackage{xcolor}
\usepackage{hyperref}
\usepackage{listings}
\usepackage{tcolorbox}
\usepackage{titlesec}

% Couleurs
\definecolor{primary}{RGB}{15, 60, 110}
\definecolor{secondary}{RGB}{41, 128, 185}
\definecolor{codebg}{RGB}{245, 245, 245}

% Configuration des titres
\titleformat{\section}{\color{primary}\Large\bfseries}{\thesection}{1em}{}[{\color{primary}\titlerule}]
\titleformat{\subsection}{\color{secondary}\large\bfseries}{\thesubsection}{1em}{}

% Configuration du code
\lstset{
    backgroundcolor=\color{codebg},
    basicstyle=\ttfamily\footnotesize,
    breaklines=true,
    frame=single,
    rulecolor=\color{lightgray},
    language=Java,
    keywordstyle=\color{blue}\bfseries,
    stringstyle=\color{red},
    commentstyle=\color{green!60!black}
}

\title{\textbf{\Huge Rapport Technique de Projet}\\[0.5cm] \Large Système de Gestion du Trafic de Fret Portuaire \\ (Port Autonome de Douala)}
\author{\textbf{Groupe de 10 Étudiants (TP308)}}
\date{\today}

\begin{document}

\maketitle

\vspace{1cm}

\section*{Présentation Générale du Projet}
L'application développée est un système de gestion portuaire destiné au \textbf{Port Autonome de Douala}. L'objectif est de numériser le suivi des marchandises arrivant par bateau. L'application gère deux types de cargaisons : les conteneurs standards et les conteneurs réfrigérés, ces derniers nécessitant une surveillance thermique stricte. 

Techniquement, le système est une application Desktop en \textbf{Java Swing} qui intègre :
\begin{itemize}
    \item Une modélisation \textbf{Orientée Objet} robuste (héritage, classes abstraites).
    \item Une interface graphique fluide et réactive avec \textbf{FlatLaf} (support mode clair/sombre).
    \item Un simulateur \textbf{Multithread} d'évolution des températures.
    \item Une \textbf{Persistance binaire} des données pour la sauvegarde de l'état.
\end{itemize}

\newpage
\section{Organisation et Répartition des Tâches}

Le groupe a fonctionné comme une équipe technique agile, divisée en 4 pôles spécialisés.

% --- EQUIPE 1 ---
\subsection{1. Équipe Core \& Métier}
\textbf{Membres :} Bahebeg, Olivia, Ali \\[0.2cm]
\textbf{Mission :} Modélisation de la POO, classes abstraites, héritage et algorithmes métier.\\[0.2cm]
L'équipe a conçu la structure de base des objets métier. Le fichier central est la classe abstraite \texttt{Marchandise.java} qui encapsule les attributs communs (ID, poids, provenance, statut). Les sous-classes \texttt{ConteneurStandard.java} et \texttt{ConteneurRefrigere.java} en héritent.
L'équipe a implémenté le \textbf{polymorphisme} sur le calcul des taxes portuaires via la méthode abstraite :
\begin{lstlisting}
// Dans Marchandise.java
public abstract double calculerTaxe();

// Dans ConteneurRefrigere.java (taxe = 5% poids + 2% surcharge energetique)
@Override
public double calculerTaxe() {
    return getPoids() * 0.05 + getPoids() * 0.02;
}
\end{lstlisting}

% --- EQUIPE 2 ---
\subsection{2. Équipe Persistance \& Données}
\textbf{Membres :} Ali Fadel, Abdel \\[0.2cm]
\textbf{Mission :} Écriture des modules d'accès aux données (Fichiers sérialisés).\\[0.2cm]
Toute la logique d'enregistrement réside dans \texttt{GestionDonnees.java}. Cette équipe a mis en place la sérialisation binaire permettant de sauvegarder la liste complète (collection \texttt{List<Marchandise>}) dans un fichier \texttt{port\_douala.ser}. Pour assurer la sécurité lors du multithreading, l'équipe a intégré un système de \textbf{verrous} (\texttt{synchronized (verrou)}).
\begin{lstlisting}
// Dans GestionDonnees.java
public void sauvegarder(String chemin) throws IOException {
    synchronized (verrou) {
        try (ObjectOutputStream oos = new ObjectOutputStream(
                new FileOutputStream(chemin))) {
            oos.writeObject(new ArrayList<>(cargaisons));
        }
    }
}
\end{lstlisting}

% --- EQUIPE 3 ---
\subsection{3. Équipe IHM Swing}
\textbf{Membres :} Alamine, Bayi, Amayen \\[0.2cm]
\textbf{Mission :} Conception des fenêtres, agencement des composants (Layouts), événements.\\[0.2cm]
L'interface graphique est orchestrée par \texttt{FenetrePrincipale.java}. L'équipe a géré l'affichage de la \texttt{JTable} avec un rendu coloré personnalisé dans \texttt{RenduTableauCargaison.java} (badges de statut arrondis). Les formulaires de saisie (\texttt{DialogueCargaison.java}) utilisent un \texttt{GridBagLayout} enveloppé dans un composant \texttt{Scrollable} pour éviter le rognage sur petit écran.

Un des éléments majeurs de l'IHM est le tableau de bord qui affiche des statistiques et le composant \texttt{GraphiqueTemperature.java}. \textbf{Ce graphique dynamique} utilise \texttt{Graphics2D} pour dessiner l'évolution en temps réel des températures des conteneurs réfrigérés à l'aide de courbes fluides et d'un quadrillage repéré.

% --- EQUIPE 4 ---
\subsection{4. Équipe Performance \& Multithreading}
\textbf{Membres :} Azefack, Lauraine \\[0.2cm]
\textbf{Mission :} Traitements asynchrones et synchronisation des données avec l'UI.\\[0.2cm]
Pour répondre à la contrainte de temps réel sans bloquer l'interface, l'équipe a implémenté la classe \texttt{MoniteurTemperature.java} qui tourne sur un thread d'arrière-plan (via l'interface \texttt{Runnable}). Ce thread se réveille toutes les 3 secondes pour modifier la température des conteneurs réfrigérés. 
Pour mettre à jour le \texttt{GraphiqueTemperature} et le tableau des alertes en toute sécurité, ils ont utilisé \texttt{SwingUtilities.invokeLater} :
\begin{lstlisting}
// Dans MoniteurTemperature.java
@Override
public void run() {
    while (actif) {
        try {
            Thread.sleep(3000); // Pause de 3 secondes
            gestionDonnees.appliquerVariationTemperature();
            
            // Mise a jour de l'IHM (JTable, Graphique) de maniere Thread-Safe
            SwingUtilities.invokeLater(() -> {
                if (onActualisation != null) onActualisation.run();
            });
        } catch (InterruptedException e) {
            Thread.currentThread().interrupt();
            break;
        }
    }
}
\end{lstlisting}

\end{document}
